% Suplemento de trabajo. No sustituye registro_k.tex ni cierra por decreto E108.
% Procedencia y alcance: INFORME_RECUPERACION.md, misma carpeta.
\subsection{Evaluación incidencial anterior al registro dodecafásico}
\label{subsec:registro-evaluacion-incidencial}

La recuperación de un registro a partir de sus diferencias cíclicas y su
carga total es una cuestión distinta de la evaluación de esas coordenadas
sobre el registro de incidencias. El primer problema se resuelve mediante el
extractor transversal de la proposición~\ref{prop:k-transversal}.
Para el segundo, una presentación histórica del
registro proporciona tres recuentos enteros en cada ventana. Es útil
explicitar su composición y determinar exactamente la información que
requiere su evaluación.

Sea \(\mathscr L\) un libro de visitas y trazas orientadas previamente
generado. Cada visita conserva la ruta, la posición APP, la hoja, el
instante, la multiplicidad y la incidencia de frontera. La evaluación
aritmética produce el entero y su descomposición
\[
 n_t=r_t+9q_t,\qquad 1\leq r_t\leq9.
\]
El entero \(n_t\) se calcula a partir de la hoja de suma o producto; su
residuo no se recibe como un valor independiente. La distinción entre
incidencia de corona e incidencia interior permanece explícita para las
visitas de residuo \(9\).

En la ventana \(E_m\), la presentación por recuentos considera
\begin{align}
 A_m&=\sum_{v\in E_m}w(v)
             \mathbf1_{\{1,3,5,7\}}(r(v)),\\
 C_m&=\sum_{j\geq0}\bigl(\nu^-_{120,m}(j)-\nu^+_{120,m}(j)\bigr),\\
 V_m&=\sum_{v\in E_m}w(v)\mathbf1_{\{9\}}(r(v))\epsilon_\partial(v).
\end{align}
Aquí \(w(v)\) es la multiplicidad incidencial y
\(\epsilon_\partial(v)=1\) en la corona, \(-1\) en el interior.
Las trazas contadas por \(\nu^\pm\) deben enumerarse mediante la regla
correspondiente del artículo. La suma entera requiere soporte finito o una
construcción alternativa que determine su significado. El número de
incidencias cronológicas y la longitud reducida \(120(1+3j)\) tienen
tipos distintos; su relación exige la unidad de longitud del lector.

Se define \([x]_{10}\in\{0,\ldots,9\}\) como el representante
euclídeo y se conservan los cocientes
\[
 x=[x]_{10}+10\lfloor x/10\rfloor.
\]
El bloque incidencial es
\begin{equation}
 z_m=100[A_m]_{10}+10[C_m]_{10}+[V_m]_{10}.
 \label{eq:registro-bloque-incidencial}
\end{equation}
Por composición se obtiene la aplicación
\begin{equation}
 \mathcal E_{\rm cnt}(\mathscr L)=
 \left((z_m-z_{m+3})_{m=1}^{12},
       (z_m-z_{m+4})_{m=1}^{12},\sum_{m=1}^{12}z_m\right).
\label{eq:registro-composicion-incidencial}
\end{equation}
Los índices se leen cíclicamente módulo \(12\). En este apartado,
\((D_jz)_m=z_m-z_{m+j}\), para \(j=3,4\), y
\(Q(z)=\sum_{m=1}^{12}z_m\); por tanto, la composición anterior es
\((D_3z,D_4z,Q(z))\). Esta convención se conserva en
\S\ref{subsec:k-transversal}, donde se demuestra la reconstrucción
integral completa.
Esta fórmula produce sus coordenadas a partir de los recuentos. No utiliza
\(K\), \(U\), \(\alpha\) ni las coordenadas transversales esperadas
como argumentos. Su identificación con
\(\mathcal E_{108}^{90,120}\) requiere componerla con la familia y
las incidencias efectivas del estado terminal; no se deduce únicamente
de la invertibilidad del sistema transversal.

\begin{proposition}[Información aritmética suficiente y necesaria]
\label{prop:registro-36-residuos}
Sean \(N,N'\in\mathbb Z^{12\times3}\), ordenados por los recuentos
\((A,C,V)\). Para la aplicación definida por
\eqref{eq:registro-bloque-incidencial}--\eqref{eq:registro-composicion-incidencial},
\[
 \mathcal E_{\rm cnt}(N)=\mathcal E_{\rm cnt}(N')
 \quad\Longleftrightarrow\quad
 N-N'\in10\mathbb Z^{36}.
\]
En particular, los \(36\) residuos sectoriales determinan exactamente
la salida de \(25\) coordenadas. La conservación de los cocientes y
de la procedencia constituye información adicional del registro.
\end{proposition}

\begin{proof}
La igualdad módulo \(10\) produce los mismos bloques y, por tanto,
las mismas coordenadas transversales. Recíprocamente, si las coordenadas
transversales coinciden, \(D_3(z-z')=D_4(z-z')=0\). Los pasos
\(3\) y \(4\) generan el ciclo de longitud \(12\), de modo que
\(z-z'\) es constante. La igualdad de cargas totales anula esa
constante. Finalmente, la representación decimal de un entero entre
\(0\) y \(999\) tiene tres dígitos únicos. Se recuperan así los
tres residuos de cada ventana. La afirmación es una identidad algebraica
para todo \(N,N'\), no una inferencia obtenida de pruebas finitas.
\end{proof}

\subsection{Conjugación de los recuentos y términos de frontera}
\label{subsec:registro-conjugacion-incidencial}

Sea \(\rho(m)=13-m\). Considérese una conjugación de incidencias
que invierte el orden de las ventanas, conserva la clase aritmética y la
clase de frontera de las visitas e intercambia los canales de las trazas.
Entonces
\[
 (A_m^\dagger,C_m^\dagger,V_m^\dagger)
 =(A_{\rho(m)},-C_{\rho(m)},V_{\rho(m)}).
\]
Si \(c_m=[C_m]_{10}\), la reducción decimal produce la identidad
\begin{equation}
 z_m^\dagger=z_{\rho(m)}
       +100\mathbf1_{\{c_{\rho(m)}\ne0\}}-20c_{\rho(m)}.
 \label{eq:registro-conjugacion-decimal}
\end{equation}
En efecto, \([-C]_{10}=0\) cuando \([C]_{10}=0\), y
\([-C]_{10}=10-[C]_{10}\) en otro caso. La fórmula conserva
exactamente esa diferencia. La negación del canal no equivale a negar
el bloque decimal completo.

La aplicación de esta identidad a una involución concreta del TPK exige
verificar su acción sobre las incidencias. Para un lector de ocupaciones
evaluado antes del avance, la inversión de una ruta
\(\gamma=(x_0,\ldots,x_n)\) satisface
\[
 \sum_{k=0}^{n-1}f(x_{n-k})
 =\sum_{k=0}^{n-1}f(x_k)+f(x_n)-f(x_0).
\]
Los términos de los extremos se incorporan antes de reducir módulo
\(10\). En rutas cerradas se anulan; en ventanas abiertas pueden
alterar sus residuos. Esta corrección, ya presente en el desarrollo
bidireccional del transporte, especifica qué debe conservar el empalme
entre la ruta y los recuentos.

\begin{proposition}[Necesidad de la información no periódica del registro]
Supóngase que el observable de una familia fija de rutas satisface
\(o(t+54)=o(t)\), que la multiplicidad se conserva bajo ese retorno
y que el mismo lector local actúa en las ventanas de \(9\) instantes.
Entonces sus recuentos y bloques satisfacen
\[
 (A,C,V)_{m+6}=(A,C,V)_m,\qquad z_{m+6}=z_m.
\]
\end{proposition}

\begin{proof}
La traslación \(t\mapsto t+54\) lleva \(E_m\) biyectivamente
a \(E_{m+6}\), conserva cada incidencia contada y su multiplicidad.
La igualdad se transmite a la suma y a la reducción decimal.
\end{proof}

Por tanto, una presentación con \(12\) bloques distintos requiere que
el lector conserve alguna información que no factorice por ese observable
de período \(54\): selección de incidencias, orientación efectiva,
extremos, multiplicidad o memoria. La proposición delimita una reducción
inadmisible del estado; no afirma que la dinámica completa tenga
período \(54\) ni determina por sí misma cuál es su lector terminal.
